\begin{table*}[t]
  \centering
  \caption{Operational failure taxonomy from the methodology asset. This definition table is evidence-tagged, not a result table; it supports the C3 design implication by naming what a console should surface.}
  \label{tab:failure-taxonomy}
  \scriptsize
  \begin{tabular}{L{0.15\linewidth}L{0.25\linewidth}L{0.22\linewidth}L{0.28\linewidth}}
    \toprule
    Class & Observable symptom & Evidence status & Deployment meaning \\
    \midrule
    F1. Legible-but-non-transfer facade & Axis is representationally legible, but steering does not beat best-prompt behavior & Supported in tested scope: E-0003 plus E-0005/E-0006 & Treat latent direction as a diagnostic probe, not production control \\
    F2. Calibration backfire under steering & Uncertainty/calibration axis worsens under steering & Supported/generalized in the tested 2$\times$2: E-0005 and E-0006; mechanism open after E-0007 & Do not deploy naive CAA/ITI steering as a trust-control primitive for calibration-facing products \\
    F3. Coherence-fragility corridor & Higher steering strength risks coherence collapse or forces conservative usable $\alpha$ & Instrumented by preregistered coherence gates; not triggered in tested cells & Stop latent route if coherence deteriorates before meaningful gain \\
    F4. Layer-sensitive instability & Axis effect depends strongly on layer choice & Supported in scope by E-0003 layer sensitivity and E-0006 per-model layer re-derivation & Require layer sweeps and stability checks before deployment \\
    F5. Axis extraction degeneration & Extracted direction overshoots or fails to encode intended construct & Supported in scope by the focus-axis overshoot/no-facade in E-0003 & Exclude the axis from control claims until redesigned \\
    \bottomrule
  \end{tabular}
\end{table*}
